%% Checkliste – erscheint nur auf der Kontrollseite im Entwurf.
%% Quellen: Portfolios im Projektpraktikum 2026-2027 (Marti, 09.07.2026),
%% offizielle Gerüste 2026-2027, Zitierregeln (31.08.2026),
%% Abbildungen und Gleichungen (Limmer 2007), Moodle-Terminliste.

\textbf{Fristen (Moodle; nächste Abgabe wird erst nach der vorherigen freigeschaltet!)}\par\smallskip
{\small\begin{tabular}{@{}ll@{}}
21.\,10.\,2026 & Motivationsformular (ohne Reflexion) -- obligatorisch\\
28.\,10.\,2026 & Arbeitsplan + Reflexion -- obligatorisch\\
11.\,11.\,2026 & Wissenschaftliche Arbeit + Reflexion -- obligatorisch\\
09.\,12.\,2026 & Protokoll eines Praktikumstages + Reflexion -- obligatorisch\\
20.\,01.\,2027 & Messergebnis + Reflexion -- obligatorisch\\
23.\,01.\,2027 & Späteste Anmeldung mündliche Leistung (Portal)\\
27.\,01.\,2027 & Folien des Vortrags + Materialien für Mitstudierende\\
28.\,01.\,2027 & Vorträge (Prüfung, Anwesenheitspflicht)\\
07.\,02.\,2027 & Folien mit Kommentar + Reflexion (vorgeschlagen, freiwillig)\\
21.\,02.\,2027 & Aussage aus anderem Vortrag + Reflexion (vorgeschlagen, freiwillig)\\
10.\,03.\,2027 & Späteste Anmeldung schriftliche Leistung (Portal)\\
\textbf{15.\,03.\,2027} & \textbf{Endabgabe Portfolio -- verspätet = nicht bestanden}\\
\end{tabular}}\par\medskip

\textbf{Vor jeder Abgabe prüfen}
\begin{itemize}\itemsep0pt
\item \texttt{\textbackslash Abgabetrue} gesetzt, keine grüne Farbe, keine Hilfen, \emph{0 Platzhalter}.
\item Deckblatt vollständig: Name, Vorname, Matrikelnummer, Studiengang, E-Mail, Thema (Projekt + Schwerpunkt).
\item Inhaltsverzeichnis mit Seitenzahlen, auch der Anhänge; durchgehend Seitenzahlen.
\item Längen eingehalten (siehe Tabelle oben; Abbildungen/Tabellen/Legenden zählen nicht mit):
      Einleitung 2--4, Beschreibung $\tfrac14$--$\tfrac12$, Diskussion der Physik 1--3,
      Reflexion 1--2, Gesamtreflexion 3--6 Seiten. Schriftgröße 12\,pt (fest eingestellt).
\item Beschreibung und Reflexion sind \emph{keine Nacherzählung}; die Leitfragen werden in Fließtext
      beantwortet (nicht in Einzelsätzen), erkennbar zugeordnet.
\item Paper (Einlage 2) wird in den späteren Einlagen zur Argumentation verwendet; Gleichungen werden
      zur Argumentation genutzt.
\item Peer-Feedback (Kommilitonen) eingeholt, aufgegriffen -- und zitiert.
\item Jede Übernahme zitiert: Literatur, Bilder (\texttt{\textbackslash cite[Abb.~3.4]\{key\}}),
      Gleichungen (\texttt{\textbackslash cite[Gl.~(3)]\{key\}}), Aussagen von Personen, Anleitungen,
      Datenblätter, Webseiten (mit Abrufdatum), \textbf{jede KI-Nutzung mit System und ALLEN Prompts}.
\item Journalartikel: Autoren, Titel, Journal, Band, Seiten, Jahr, DOI (nicht als Link).
\item Abbildungen: Bildunterschrift, im Text referenziert (\texttt{\textbackslash ref}), Schrift
      $\geq 10$\,pt, Achsen mit Einheit als \enquote{$f$ / kHz} (nicht \enquote{[kHz]}), ggf.\ Residuenplot.
\item Tabellen: Überschrift, im Text referenziert; Messdaten nicht als lange Listen im Text, sondern im Anhang.
\item Wichtige Gleichungen nummeriert, alle Symbole definiert und einheitlich verwendet.
\item Dateiname für Moodle: \texttt{Kautzsch\_WelcheEinlage.pdf} (nur PDF oder DOCX).
\item Bis auf den Arbeitsplan müssen sich alle Einlagen von denen der Partnerin/des Partners unterscheiden.
\end{itemize}
